\chapter{Statistická analýza výsledků a škálování}
\label{chap:vysledky}

V této kapitole prezentujeme empirické výsledky faktorového experimentu, dekompozici rozptylu (ANOVA), lineární smíšené modely a odhadnuté parametry mocninných škálovacích zákonů.

\section{Naměřené hodnoty BPC a gramatické přesnosti}
Tabulka \ref{tab:vysledky_grid} shrnuje průměrné hodnoty validačního BPC a úspěšnosti v gramatickém testu pro jednotlivé konfigurace.

\begin{table}[htbp]
    \centering
    \caption{Výsledky faktorového pokusu na datech toki pona}
    \label{tab:vysledky_grid}
    \begin{tabular}{llcccc}
        \toprule
        Tokenizér & Model & Frakce dat ($D$) & Val BPC $\downarrow$ & Znaková PPL $\downarrow$ & Gramatická přesnost $\uparrow$ \\
        \midrule
        CHAR & micro & 25\,\% & 1{,}8495 & 3{,}60 & 39{,}0\,\% \\
        CHAR & micro & 75\,\% & 1{,}4178 & 2{,}67 & 53{,}7\,\% \\
        CHAR & mini  & 25\,\% & 1{,}3867 & 2{,}61 & 58{,}5\,\% \\
        CHAR & mini  & 75\,\% & \textbf{1{,}0728} & \textbf{2{,}10} & 85{,}4\,\% \\
        \midrule
        BPE  & micro & 25\,\% & 1{,}6536 & 3{,}15 & 75{,}6\,\% \\
        BPE  & micro & 75\,\% & 1{,}4126 & 2{,}66 & 85{,}4\,\% \\
        BPE  & mini  & 25\,\% & 1{,}4079 & 2{,}65 & 82{,}9\,\% \\
        BPE  & mini  & 75\,\% & 1{,}2183 & 2{,}33 & 92{,}7\,\% \\
        \midrule
        WORD & micro & 25\,\% & 1{,}2848 & 2{,}44 & 85{,}4\,\% \\
        WORD & micro & 75\,\% & 1{,}1933 & 2{,}29 & 87{,}8\,\% \\
        WORD & mini  & 25\,\% & 1{,}1987 & 2{,}30 & 90{,}2\,\% \\
        WORD & mini  & 75\,\% & \textbf{1{,}1185} & \textbf{2{,}17} & \textbf{100{,}0\,\%} \\
        \bottomrule
    \end{tabular}
\end{table}

\section{Analýza rozptylu (ANOVA) a velikost účinku (\texorpdfstring{$\eta^2$}{eta2})}
Pro ověření statistické významnosti vlivu faktorů byl aplikován model ANOVA typu II. Celkový model vysvětluje 99{,}25\,\% rozptylu BPC:
\begin{itemize}
    \item \textbf{Model Size ($N$):} $F = 119{,}21$, $p = 0{,}0016$, $\eta^2 = 29{,}83\,\%$,
    \item \textbf{Data Size ($D$):} $F = 109{,}14$, $p = 0{,}0019$, $\eta^2 = 27{,}31\,\%$,
    \item \textbf{Tokenizer ($T$):} $F = 50{,}27$, $p = 0{,}0049$, $\eta^2 = 25{,}16\,\%$,
    \item \textbf{Interakce Model $\times$ Tokenizer:} $F = 18{,}97$, $p = 0{,}0198$, $\eta^2 = 9{,}50\,\%$,
    \item \textbf{Interakce Data $\times$ Tokenizer:} $F = 14{,}89$, $p = 0{,}0277$, $\eta^2 = 7{,}45\,\%$,
    \item \textbf{Reziduální rozptyl:} pouhých $0{,}75\,\%$.
\end{itemize}

Všechny hlavní efekty i dvojné interakce jsou statisticky signifikantní na hladině významnosti $\alpha = 0{,}05$.

\section{Odhad empirických škálovacích zákonů}
Nelineární regrese metody nejmenších čtverců odhadla parametry rovnice \eqref{eq:scaling_law} na hodnoty:
\begin{equation}
    \text{BPC}(N, D) = 0{,}555 + \frac{77{,}05}{N^{0{,}462}} + \frac{93{,}78}{D^{0{,}390}}.
\end{equation}

Odhadnutá hodnota $E \approx 0{,}555$ bitů na znak představuje empirický odhad entropie jazyka toki pona. Exponent $\alpha = 0{,}462$ a $\beta = 0{,}390$ potvrzují, že v malém měřítku přináší zvětšování kapacity modelu ještě výraznější relativní přínos než v asymptotické oblasti velkých modelů.

\begin{figure}[htbp]
    \centering
    \includegraphics[width=0.85\textwidth]{scaling_data_bpc.pdf}
    \caption{Škálování validačního BPC v závislosti na objemu dat pro jednotlivé tokenizéry.}
    \label{fig:scaling_data}
\end{figure}

\begin{figure}[htbp]
    \centering
    \includegraphics[width=0.85\textwidth]{interaction_effects.pdf}
    \caption{Interakční graf mezi velikostí korpusu a typem tokenizéru.}
    \label{fig:interaction}
\end{figure}

\begin{figure}[htbp]
    \centering
    \includegraphics[width=0.85\textwidth]{grammar_vs_bpc.pdf}
    \caption{Závislost gramatické přijatelnosti (BLiMP) na Bits-Per-Character.}
    \label{fig:grammar_vs_bpc}
\end{figure}
